\section{Introduction}
\label{sec:introduction-v132}

Which defining equations can be recovered from an abstract degeneracy
scheme?  Reduced rank loci often cannot distinguish the equations:
large families of presentations have the same determinantal support.
We prove that the nonreduced multiplication-failure scheme does
recover the quadratic relations in two natural families: every
basepoint-free four-dimensional web with smooth Jacobian quartic,
and every pencil of quadrics in every dimension at least three.
The pencil theorem has no regularity hypothesis.  In both cases the
reduced scheme alone is a fixed Schubert divisor.

The geometric mechanism is intrinsic coefficient extraction.  The
scheme determines its deepest stratum, the degree-one bundle of the
normal cone, and an oriented rank-one ruling.  An ideal quotient by
the determinant line recovers coefficient subspaces, all transported
by one projective coordinate transformation.  We prove directly from
the actual multiplication matrix that this construction applies to
a whole Grassmannian family of quadratic relation spaces.  For webs,
an exterior-support theorem uniquely recombines two Pluecker
components.  For pencils the relevant representation is irreducible,
so the entire Pluecker line is recovered at once.  This last result
distinguishes pencils with the same discriminant, including its
multiplicities, but different elementary divisors.  It also gives a
fixed-reduction scheme invariant of hyperelliptic curves through the
classical correspondence with simple-discriminant pencils.

The first nilpotent layer and the deeper boundary retain additional
geometry: a polarized Jacobian K3 surface, contact divisors on a Segre
quadric, vertex-primary corrections, and their higher-rank
specializations.  These are intrinsic layers, not artifacts of the
presentation used to compute them.  Their complete theory is retained
in the technical supplement, separately from the proof of the main
inverse theorem.

\subsection{The multiplication failure scheme}

Let \(V\) be a complex vector space of dimension
\(e\in\{3,4\}\), put \(p=e(e-1)/2\), and let
\[
 R\subset\Sym^2V
\]
be an \(e\)-dimensional relation space.  Set
\[
 S_R=\Sym^2V/R,\qquad
 B_R=\C\oplus V\oplus S_R,
\]
with multiplication induced by the quotient
\(\Sym^2V\twoheadrightarrow S_R\) and with
\((V\oplus S_R)S_R=0\).  On
\[
 X_R=\Gr(e,V\oplus S_R)
\]
let \(\mathcal K\) be the tautological subbundle.  For every
\(m\ge2\) the multiplication map defines the same zeroth-Fitting
scheme
\begin{equation}\label{eq:failure-intro}
 \widehat D_R=
 V\!\left(
 \Fitt_0\coker[
 \Sym^m(\OO\oplus\mathcal K)\longrightarrow
 B_R\otimes\OO]\right).
\end{equation}
Its reduction is the Schubert divisor
\(\widehat\Delta_R\) on which \(\mathcal K\to V\) drops rank.

Regard \(R\) as a web of quadrics on \(\Pj(V^*)\).  On the
invariant open \(G_e^\circ\) used throughout, the web is
basepoint-free and its Jacobian hypersurface
\begin{equation}\label{eq:Y-intro}
 Y_R=V\!\left(\det(\partial q_i/\partial x_j)\right)
 \subset\Pj(V^*)
\end{equation}
is smooth.  For \(e=4\) it is a quartic K3 surface.


\subsection{Recovering the original web}

The central inverse theorem uses the full nonreduced scheme,
including information at the deepest projection-corank stratum.  The
following conclusion holds without an additional genericity condition
and is proved in
Section~\ref{sec:intrinsic-web-reconstruction}.

\begin{theorem}[Intrinsic web reconstruction]
\label{thm:main-web-reconstruction}
For every pair \(R,R'\in G_4^\circ\),
\[
 \widehat D_R\simeq\widehat D_{R'}\quad\Longrightarrow\quad
 R'=gR\ \text{for some }g\in\PGL(V).
\]
Consequently the abstract failure scheme determines \(B_R\)
throughout the smooth Jacobian locus and separates every finite
web-to-K3 packet contained in this locus.
\end{theorem}

The proof first recovers \(\Gr(4,S_R)\) by iterated reduced
singular loci.  Its nonreduced normal cone identifies the two
components of
\[
 \bigwedge^4\Sym^2V=
 \mathbb S_{(4,3,1,0)}V\oplus\mathbb S_{(5,1,1,1)}V.
\]
The second component is the Jacobian covariant; the first is not
retained by the quartic alone.  The two projective components define a
line in the Pluecker space.  The contraction-kernel theorem
(Lemma~\ref{lem:schur-contraction-kernel}) computes the exact exterior
support of the Jacobian component.  For a smooth quartic that support
has dimension ten.  A secant or tangent four-vector on the Grassmannian
has support at most eight.  This excludes every second point and every
infinitesimal ambiguity, so the line meets the Grassmannian in one
reduced point.  In fact three essential variables already suffice.  The global tensor rulings of the normal cone
exclude the transposition ambiguity of an isolated square-matrix
cone.  No extension of an abstract scheme isomorphism to the ambient
Grassmannian is assumed.


The exceptional fibres of the two-component invariant are described in
Theorem~\ref{thm:exceptional-component-fibres}.  A proper intersection
of the component line with the Grassmannian has length at most two.
Its second point can coincide with the first, lie at a discarded pure
endpoint, or give a distinct candidate exchanged by a regular residual
involution.  A positive-dimensional fibre is exactly a Grassmann line
specified by a three-plane inside a five-plane.  The tangent condition
\(Qw_R\in\bigwedge^3R\wedge W\) distinguishes the double-point
case from a reduced secant pair.  All these alternatives are excluded on \(G_4^\circ\) by
Theorem~\ref{thm:global-schur-recombination}.  The ambient obstruction
is confined to the binary-Jacobian boundary, defined invariantly by the
first catalecticant.  Thus the smooth family has no residual secant
involution, flag-line ambiguity, or ramification divisor to resolve.
Outside that family the component-fibre classification is not promoted
to a classification of abstract failure-scheme isomorphisms.


\subsection{The mechanism in arbitrary dimension}

There is a second inverse theorem, for a family that is not confined
to the two-summand plethysm in dimension four.  If \(\dim V=n\),
polarization and the Jacobian determinant give natural maps
\[
 j_n:\det V\otimes\Sym^nV\longrightarrow\bigwedge^n\Sym^2V,
 \qquad C_n:\bigwedge^n\Sym^2V\longrightarrow\det V\otimes\Sym^nV,
 \qquad C_nj_n=2\,\mathrm{id}.
\]
Put \(Q_n=\frac12j_nC_n\) and \(P_n=1-Q_n\).
Theorems~\ref{thm:uniform-contraction}--\ref{thm:uniform-recombination}
prove the exact all-ranks contraction kernel, the resulting support
formula, and the following consequence.  For every \(n\ge4\), on
the nonempty open set of relation spaces whose Jacobians have at
least three essential variables, the pair
\(([P_nw_R],[Q_nw_R])\) determines \(R\).  The corresponding
pencil intersection is a single section even after nonreduced base
change.  The complement \(\ker C_n\) need not be irreducible.

The numerical mechanism is uniform: for \(d<n\) essential variables,
the exterior support is
\(\binom{n+1}{2}-\binom{n-d+1}{2}\), and for \(d=n\) it is full
except for a one-dimensional annihilator of a nondegenerate quadratic
power when \(n\) is even.  The resulting gap above \(2n\) excludes
secants and tangents.  The skew-flattening support criterion itself
is classical \cite{LandsbergOttaviani,Sheridan}; the calculation here
is its restriction to the specified Jacobian component, including
all singular and nondegenerate contraction kernels.  In dimension
four the intrinsic failure scheme supplies the required component
pair, thereby turning the uniform inverse mechanism into
Theorem~\ref{thm:main-web-reconstruction}.  The smooth application
has four essential variables; the algebraic exclusion is stronger
and already applies with three.

\subsection{Natural pencil reconstruction and discriminant fibres}

Let \(n\ge3\), \(N=n(n+1)/2\), and let
\(R\subset\Sym^2V\) now have dimension two.  Apply the same
multiplication construction to the algebra with Hilbert function
\((1,n,N-2)\).  Theorem~\ref{thm:natural-pencil-reconstruction}
proves
\[
 \widehat D_R\simeq\widehat D_{R'}\quad\Longleftrightarrow\quad
                  R'=gR\quad\text{for some }g\in\PGL(V).
\]
Theorem~\ref{thm:automatic-quadratic-coefficients} derives every
hypothesis of the coefficient criterion for the wider family
\(R\in\Gr(r,\Sym^2V)\), \(1\le r<N-n\).  Thus the
pencil result is an application of the original multiplication
failure scheme, not of an auxiliary map carrying a coefficient block.

For regular pencils, the Weierstrass--Segre classification uses
root-labelled elementary divisors, not just the determinant
\cite[Theorem~1.1 and Corollary~2.1]{FevolaMandelshtamSturmfels}.
Example~\ref{ex:same-discriminant-distinct-failure} exhibits two
pencils with identical discriminants and distinct partitions over
the unique double root.  Their failure schemes are nonisomorphic,
although their reductions and discriminant double covers agree.
For \(n=2g+2\), \(g\ge2\), the simple-discriminant locus
recovers the classical hyperelliptic moduli correspondence
\cite[\S4.2]{SmithQuadrics};
Corollary~\ref{cor:hyperelliptic-failure-invariant} gives an
injective scheme invariant of the resulting \((2g-1)\)-dimensional
family of curve classes.  The new assertion is recovery from the
unmarked failure scheme, not a new classification of pencils or curves.

For webs, the polarized K3 construction has a nine-dimensional image
and is generically finite (Section~\ref{sec:moduli}).  We do not
determine its generic degree or assume it exceeds one.  Our web
theorem separates any fibre on the stated smooth locus.  The actual
loss of information by a classical invariant, used here to illustrate
the strength of nonreduced reconstruction, is the explicit pencil
discriminant example rather than a conjectural nontrivial K3 packet.


\subsection{The first infinitesimal order that remembers the pencil}

There is an exact bound on the amount of nonreduced structure needed
for pencil reconstruction.  Put \(N=\binom{n+1}{2}\),
\(p=N-2\), and \(d=n+2p=n^2+2n-4\).  If
\(\mathfrak b\) is the ideal of the intrinsic socle Grassmannian
\(B=\Gr(n,S_R)\) in \(\widehat D_R\), let
\(Z_R^{[k]}\) denote its neighbourhood defined by
\(\mathfrak b^{k+1}\).  Theorem~\ref{thm:sharp-finite-pencil}
proves that every \(Z_R^{[k]}\) with \(k<d\) is independent
of \(R\), whereas the abstract scheme \(Z_R^{[d]}\)
determines \(R\) up to projective equivalence.  The first
nonzero homogeneous relation in its intrinsic conormal algebra
recovers the entire normal cone.  Thus the result specifies a sharp
order, rather than merely asserting that some nilpotent structure
contains additional information.

Theorem~\ref{thm:finite-parameter-immersion} strengthens the
coefficient calculation at the level of parameter schemes: in fixed
tensor coordinates the degree-\(d\) relation spaces define a
closed immersion of the whole \(\Gr(2,\Sym^2V)\).
It detects infinitesimal variations on singular as well as regular
pencils, with no auxiliary coefficient block added to the original
multiplication map.  For hyperelliptic curves of genus \(g\),
Corollary~\ref{cor:finite-hyperelliptic} gives the exact order
\(4g^2+12g+4\).  The intrinsic assertion uses a neighbourhood
of the entire socle Grassmannian; the parameter-family assertion
retains its tensor coordinates.  Neither asserts an equivalence
of unmarked moduli stacks.

\subsection{Further coefficient constructions}

For \(\dim V=n\ge2\), \(d\ge1\), and
\(0\ne L\subset\Sym^dV^*\), consider \(B=\Gr(n,H)\),
where \(\dim H>n\).  On \(\Tot\Hom(\mathcal U,V)\), the
block map \(\det T\oplus\operatorname{ev}_{L,T}\) combines
the determinant with the coefficients of the pulled-back system.
Theorem~\ref{thm:polar-system-torelli} proves that its zeroth-Fitting
scheme \(Z_L\) determines \(L\) up to projective equivalence,
from an abstract isomorphism of schemes.  Every reduction is the same
determinant hypersurface.  There is no smoothness restriction on
\(L\), and no supplied marking of the zero section or tensor factors.
The structural criterion of
Theorem~\ref{thm:structural-coefficient-reconstruction} accounts for
this family as well as the quadratic multiplication families.
The polar construction is useful independently of recombination, but
the natural pencil theorem, not this auxiliary block map, supplies
the application to classical moduli in the preceding subsection.

\subsection{Intrinsic reconstruction}

On the projection-corank-at-most-one locus let \(D_R\) be the
restriction of \(\widehat D_R\), let \(\NN_R\) be its
nilradical, and let \(E_R\) be the annihilator scheme of the first
nilpotent layer.  The first-layer reconstruction is the following result
proved in Section~\ref{sec:polarization}.

\begin{theorem}[Polarized reconstruction]
\label{thm:main-reconstruction}
For \(R\in G_e^\circ\), put
\[
 \LL_R=\omega_{E_R}\otimes
 \NN_R^{\otimes-(p+e-1)}.
\]
Then
\[
 \bigoplus_{n\ge0}H^0(E_R,\LL_R^{\otimes n})
 \simeq
 \bigoplus_{n\ge0}H^0(Y_R,\OO_{Y_R}(en))
\]
as graded algebras.  In particular, for \(e=4\) the abstract
failure scheme determines the polarized quartic K3
\((Y_R,\OO_{Y_R}(1))\).
\end{theorem}


\subsection{Organization and relation to prior work}

Section~\ref{sec:uniform-support} proves the dimension-uniform
contraction and component-pair inverse theorems.
Sections~\ref{sec:relations} and \ref{sec:fitting} construct the algebra
and its Fitting scheme.  Section~\ref{sec:intrinsic-web-reconstruction}
contains the complete reconstruction argument: the deepest stratum,
its oriented normal cone, the universal coefficient readout, one
common projective transformation, and the contraction-kernel and
exterior-support theorems.  The scheme structure of the cone is
used directly; no primary decomposition of its ideal is needed.
Section~\ref{sec:natural-pencils} proves automatic coefficient extraction,
the all-pencil theorem, and the discriminant and hyperelliptic
consequences.  Section~\ref{sec:finite-neighbourhoods} proves the
sharp finite-order and scheme-theoretic parameter results.  Section~\ref{sec:structural-coefficient-principle} proves the abstract
coefficient criterion and the polar-restriction theorem.
Sections~\ref{sec:polarization}--\ref{sec:classical} compare this
invariant with the polarized Jacobian surface and its moduli.
Section~\ref{sec:support-literature} identifies the classical
representation and secant tools and the precise additional
contraction calculation.

The technical supplement is a separate reading object, not part of
the proof of the main inverse theorem.  It contains the full
nilpotent-depth, primary-boundary, collision, relative-filtration,
and exact-coordinate results, as well as the detailed Reye comparison
and the ambient component-pencil classification.  Its rank-defect
schemes are defined by actual minors; their support is only contained
in the binary-Jacobian boundary.  No complete embedded-primary atlas
of the highest-corank residual layers is used or asserted.

Finite exact checks accompany the calculations, but do not formally
verify the written structural proofs.  The historically relevant
Ballico paper \cite{Ballico93} remains a specific documentary issue:
its publisher first page has been inspected, but its complete text
has not been obtained.  No theorem-level nonanticipation conclusion
is drawn from that page, its title, or later citations.
This does not change the hypotheses or conclusion of the reconstruction
theorem proved here.
